% Build with: xelatex PARC2026_track1_report.tex (twice)
\documentclass[10pt,a4paper]{article}
\usepackage[margin=17mm]{geometry}
\usepackage{fontspec}
\usepackage{xeCJK}
\usepackage{booktabs}
\usepackage{tabularx}
\usepackage{array}
\usepackage{enumitem}
\usepackage[table]{xcolor}
\usepackage[hidelinks]{hyperref}
\usepackage{titlesec}

% MS P Gothic is used when installed. The build environment falls back to
% IPAexGothic, whose proportions are closer to MS P Gothic than Noto Sans.
\IfFontExistsTF{MS PGothic}{
  \setmainfont[AutoFakeBold=2]{MS PGothic}
  \setCJKmainfont[AutoFakeBold=2]{MS PGothic}
}{
  \setmainfont[AutoFakeBold=2]{IPAexGothic}
  \setCJKmainfont[AutoFakeBold=2]{IPAexGothic}
}
\setCJKmonofont{Noto Sans Mono CJK JP}
\setsansfont{Noto Sans CJK JP}
\setmonofont{DejaVu Sans Mono}
\setlength{\parindent}{0pt}
\setlength{\parskip}{0.35em}
\setlist[itemize]{leftmargin=1.5em,itemsep=0.18em,topsep=0.2em}
\setlist[enumerate]{leftmargin=1.7em,itemsep=0.18em,topsep=0.2em}
\titleformat{\section}{\large\bfseries\color{blue!45!black}}{\thesection.}{0.5em}{}
\titleformat{\subsection}{\normalsize\bfseries\color{blue!35!black}}{}{0pt}{}
\titlespacing*{\section}{0pt}{0.7em}{0.3em}
\titlespacing*{\subsection}{0pt}{0.45em}{0.15em}
\pagestyle{plain}

\newcolumntype{L}{>{\columncolor{blue!4}\centering\arraybackslash}m{0.27\textwidth}}
\newcolumntype{Y}{>{\raggedright\arraybackslash}X}
\newcommand{\code}[1]{\texttt{\small #1}}
% Use \br inside a table cell. Unlike \\, it does not terminate the table row.
\newcommand{\br}{\newline}
\renewcommand{\arraystretch}{1.12}
\newcommand{\reportrow}[2]{\textbf{#1} & #2 \\}
\newcommand{\tablehead}{\rowcolor{blue!12}\textbf{項目} & \textbf{内容} \\}

\begin{document}

\begin{center}
  {\LARGE\bfseries PARC 2026 予選 Track 1 技術レポート}\par
  \vspace{0.35em}
  {\small 氏名：山本湧也\quad 提出物：\code{soup\_track1full\_w\_pi05.zip}\quad 作成日：2026年8月15日}
\end{center}

\section{学習・推論時の工夫点および試行錯誤}

本提出物はTrack 1の公開4タスクを単一ポリシーで処理するpi0.5-LIBEROのLoRAモデルである。

\subsection{モデル選定}

はじめに事前学習モデルとしてSmolVLAとTurboVLAを検討した。SmolVLAは軽量でLoRA学習も
容易だったがTrack 1の強い視覚摂動では十分な成功率を得られず実際のロールアウトを確認しても追加学習だけで大幅に改善することは難しいと判断した\cite{smolvla}。次にLIBEROで高い成功率と高速推論を報告するTurboVLAを検討した\cite{turbovla}。TurboVLA自体は約0.2Bパラメータと小型だが公開チェックポイントはLIBEROのsuite別であった。単一方策として統合するには追加学習や複数expertの切替機構が必要となり計算資源と提出条件の両面で採用が難しかった。最終的に異種データのco-trainingによる汎化性能を持つpi0.5\cite{pi05}のLIBERO用チェックポイントを評価したところtrack1の公開4タスクで26/32を達成した。全suiteを単一重みで扱える点も踏まえpi0.5-LIBEROを採用した。
なおpi0.5-LIBEROのコードとチェックポイントはPhysical Intelligenceが公開するopenpiに含まれているがPyTorch版とJAX版が存在する。この二つは性能が異なっておりJAX版の方がtrack1の公開4タスクで高い成功率を示したためJAX版を採用した\cite{openpi}。

\subsection{学習}

利用可能なデータ量と計算資源を考慮しスクラッチ学習や全パラメータのfine-tuningではなく事前学習済みpi0.5-LIBEROに対するLoRA fine-tuningを採用した。これにより事前学習で得られた汎化性能を維持しながらTrack 1の視覚摂動へ適応させた。

\begin{itemize}
  \item \textbf{データ変換：}
    LIBERO-plusのデータは複数episodeをparquet chunkに格納した形式で提供されている。これをJAX版openpiで学習できるようepisode単位にデコードしLeRobotDatasetとして再構築した。その過程でchunk境界のindex処理とArrow schema metadataの欠落も修正した。

  \item \textbf{学習データとデータ拡張：}
    初期学習モデルを評価したところTrack 1の中でもtomato-sauceタスクで特に低い成功率を示した。そこでこのタスクを分析したところ背景の色彩と明度に強い摂動が加えられていたため暗色化や色調変換によるaugmentationを検証した。しかしaugmentationなしの成功率43.8\%に対し適用時は6.2--31.2\%となりかえって性能が低下した。

    またtrack1の4タスクだけではなく40タスクへ学習対象を広げ12,000--20,000 stepsの比較学習を行った。LIBERO-plus上のvalidation成功率は向上したもののOmniCampus上の評価スコアは低下した（狭域レシピの0.391に対し広域データ／augmentation適用版は0.260--0.313）。詳細に分析するとタスク数の増加によって公開4タスクの学習信号が希釈されるとともにaugmentationを強めるほどaction jerkが増加していた。この結果から評価では成功率だけでなく衝突を避けて安定した軌道を生成することも重要だと判断した。最終学習には公開4タスクの実データ530 episodeのみを使用しaugmentationは適用しなかった。

  \item \textbf{軌道の滑らかさ：}
    タスクを増やしたときにスコアが下がった現象について成功率だけではなくaction jerkが重要な評価指標になっていると考え予測actionのjerkを抑える補助損失を導入した。しかし成功率と衝突率は改善しなかった。また複数のaction chunkを時間方向に平均する方式も検討したがjerkが4.59から3.47へ低下した一方成功率は85.9\%から81.2\%へ
    衝突率は3.1\%から14.1\%へ悪化し実行時間も増加したため採用しなかった。

  \item \textbf{最終的な学習方式：}
    同一条件の学習でも評価結果に最大約11ポイントの変動が確認された。
    そこで単一チェックポイントへの依存を抑えるためstep 700・750・775の
    LoRAパラメータを2:1:1で平均したweighted model soupを提出モデルとした。
    このweighted soup自体のローカル16 episode/task評価は成功率85.9\%
    （drawer 93.8\%・tomato-sauce 56.3\%・milk 100\%・stove 93.8\%）、
    衝突率3.1\%（本プロジェクト中で最良）、平均cartesian jerk 4.59であり、
    単一checkpoint（step700単体、本番0.391）と同系統の低jerkな学習レジームを
    維持しつつ、平均化によりラン間のばらつきを抑えたものである。
\end{itemize}

\subsection{推論}

推論では1回の推論で10ステップのaction chunkを生成しそのうち先頭5ステップを
実行してから再推論するreceding-horizon方式を採用した。全10ステップを実行する場合より
観測変化へ早く追従でき毎ステップ推論する場合より計算時間を抑えられる。

入力画像には学習時と同じ180度回転とpadding付きresizeを適用した。ロボット状態の
手先姿勢はEuler角ではなくaxis-angleへ変換し学習時の前処理と一致させた。
またJAXのJITコンパイルをサーバー起動時に完了させることで評価中の各推論を
10秒以内に収めた。出力actionは有限値であることを確認したうえで
$[-1,1]$へclipしている。

\clearpage
\small
\section{モデル情報}

\begin{tabularx}{\textwidth}{@{}LY@{}}
\toprule
\tablehead
\midrule
\reportrow{モデル名およびバージョン}{
  pi0.5-LIBERO LoRA LIBERO-plus 4タスク単一方策モデル}
\reportrow{ベースモデル名}{
  Physical Intelligenceの公式pi0.5-LIBEROチェックポイント(JAX版)}
\reportrow{ベース重みの入手元}{
  \url{gs://openpi-assets/checkpoints/pi05_libero}}
\reportrow{revision／commit hash}{
  openpi: GitHub \code{Physical-Intelligence/openpi}, commit：
  \code{15a9616a00943ada6c20a0f158e3adb39df2ccac}
  元データ\code{lerobot/libero\_plus}, commit：
  \code{f3f49f426d75030177b18778374005bc12ccd588}。}
\reportrow{提出チェックポイントのハッシュ値}{
  提出zipのSHA-256：
  {\scriptsize\ttfamily b98f61fc2b0349b20a4ca403f84247cb9f71a77f715d5f1ef491b147cf227241}}
\reportrow{モデル構成}{
  pi0.5: PaliGemma-2B系VLMおよび約300Mパラメータのflow-matching action expert（dtypeはbfloat16）}
\reportrow{推論方式}{
  openpi/JAXで単一チェックポイントをロードしFastAPIの\code{/health}, \code{/reset}, \code{/act}から提供する}
\reportrow{Action chunk}{
  \code{action\_horizon=10}で10ステップを予測し先頭5ステップをオープンループ実行後に再推論するreceding-horizon方式（\code{PI05\_ACTION\_CHUNK=5}）}
\bottomrule
\end{tabularx}

\section{学習情報}
\begin{tabularx}{\textwidth}{@{}LY@{}}
\toprule
\tablehead
\midrule
\reportrow{使用した学習データ}{
  HF Hubの\code{lerobot/libero\_plus}（LIBERO-plus）をベースに使用する
  独自の2段階パイプラインでopenpi形式へ変換した。公開4タスクの実演のみを使用し独自収集データ・生成データ・最終学習時の合成augmentationはない}
\reportrow{件数・タスク数・エピソード数}{
  530エピソード／74,404フレーム／4タスク／20 fps／画像256$\times$256。最終的な内訳はdrawer 60・milk 60・stove 60・tomato-sauce 350エピソードとした。}
\reportrow{学習方法}{
  pi0.5-LIBEROへのflow-matching模倣学習によるLoRAファインチューニングを行った。
  optimizerはAdamW。勾配ノルム上限は1.0。学習率はcosine decayでスケジューリングする。}
\reportrow{学習ステップ数}{
  比較実験は12,000--20,000 steps。
  提出モデルの最終レシピは1,000 stepsで成功率を安定化させるためstep 700・750・775のcheckpointの重みを2:1:1で平均。}
\reportrow{学習対象パラメータ}{
  VLMおよびaction expert内のLoRA行列のみ（約50.0M／全体の約1.47\%）。denseパラメータはすべて凍結する。}
\reportrow{区分}{
  LoRA。Full fine-tuningおよびLoRA以外のadapterは使用していない。}
\reportrow{主なハイパーパラメータ}{
  \code{peak\_lr=1e-5}／\code{decay\_lr=1e-6}／
  \code{warmup\_steps=100}／\code{batch\_size=8}／
  \code{clip\_gradient\_norm=1.0}／augmentationなし／jerk補助損失0／EMAなし。}
\bottomrule
\end{tabularx}

\section{権利関係}

\begin{tabularx}{\textwidth}{@{}LY@{}}
\toprule
\tablehead
\midrule
\reportrow{ベースモデルのライセンス}{
  openpiのコードはApache License 2.0である。
  PaliGemma／Gemma由来の重みにはこれに加えてGoogle Gemma Terms of Use（\code{LICENSE\_GEMMA.txt}）が適用され利用制限および再配布時の告知義務がある。}
\reportrow{学習データのライセンス}{
  LIBERO-plusには確認時点で明示的なLICENSEおよびREADME上のライセンス記載がない。基盤のLIBERO由来部分はMIT License（Copyright (c) 2023 Lifelong Robot Learning）である。}
\reportrow{提出物内の第三者コード}{
  openpiコードは上記の通り。JAX/Flax・FastAPI・NumPy等の依存は主にApache-2.0・BSD・MIT系である。
  詳細は\code{THIRD\_PARTY\_LICENSES.md}および提出zip内の\code{OPENPI\_LICENSE}, \code{LEROBOT\_LICENSE}を参照。}
\bottomrule
\end{tabularx}

\renewcommand{\refname}{参考文献}
{\footnotesize
\begin{thebibliography}{9}
\setlength{\itemsep}{0pt}
\bibitem{smolvla} M. Shukor et al., ``SmolVLA: A Vision-Language-Action Model for Affordable and Efficient Robotics,'' arXiv:2506.01844, 2025.
\bibitem{turbovla} H. Xie et al., ``TurboVLA: Real-Time Vision-Language-Action Model at 32 Hz on an RTX 4090 with $<$1 GB VRAM,'' arXiv:2607.27205, 2026.
\bibitem{pi05} Physical Intelligence et al., ``$\pi_{0.5}$: a Vision-Language-Action Model with Open-World Generalization,'' arXiv:2504.16054, 2025.
\bibitem{openpi} Physical Intelligence, ``openpi,'' \url{https://github.com/Physical-Intelligence/openpi}.
\end{thebibliography}
}

\end{document}
